\section*{9. Conclusion}
PipeViz demonstrates how cycle-accurate pipeline simulation, compiler experimentation, and LLM-assisted reasoning can be combined into a unified platform for micro architectural analysis. By connecting source code, generated AArch64 assembly, hazard detection, and pipeline visualization within a single interactive workflow, the system enables users to better understand the relationship between compiler optimizations and processor execution behavior.

The integration of an LLM assistant further extends the platform beyond traditional simulators by enabling context-aware explanations and optimization recommendations grounded in actual execution traces. Through the presented case studies, PipeViz shows how techniques such as loop unrolling and cache blocking influence hazard frequency, instruction-level parallelism, and pipeline efficiency across different scheduling architectures.

Overall, PipeViz serves both as an educational framework for teaching computer architecture concepts and as an experimental environment for exploring optimization strategies in modern processor pipelines. The project highlights the potential of combining AI-assisted reasoning with low-level architectural simulation to create more interactive and insightful performance analysis tools.